\documentclass[aps,pre,twocolumn,superscriptaddress,floatfix,longbibliography]{revtex4-2}

\usepackage{amsmath}
\usepackage{amssymb}
\usepackage{graphicx}
\usepackage{bm}
\usepackage{xcolor}

\begin{document}

\section{Model}

\begin{figure*}[!ht]
    \centering
    \includegraphics[width=\linewidth]{fig2/figure_composite.pdf}
    \caption{\label{fig:env}Pipeline for estimating the mutual information of the array.
    \textbf{(A)} Two-dimensional UMAP projection of the $D$-dimensional
    morphochemical latent space. Ligands (dots) are colored by chemical family,
    receptors (open circles) are numbered. The affinity of a receptor for a
    ligand is set by their distance $d = \|\mathbf{v}_u - \mathbf{v}_\ell\|$
    (red arrow).
    \textbf{(B)} Activation threshold $c^*$ as a function of $d$, saturating
    between $c^*_\text{min}$ at complete match and $c^*_\text{max}$ at complete
    mismatch.
    \textbf{(C)} Monte Carlo loop. The environment draws the per-receptor
    threshold vector $\mathbf{c}^*_{\mathbf{r},\ell}$ over the $L$ ligands and a
    mixture $\mathbf{c} = (n_1 c_1,\dots,n_L c_L)$, with presence
    $n_\ell \sim \text{Bernoulli}(p)$ and concentration
    $c_\ell \sim \text{LogNormal}(\mu,\sigma)$. The physics step turns this into
    a binary activation pattern $A = (1,0,\dots,1)$. Repeating over sensing
    events builds the empirical distribution $p(A)$, from which the mutual
    information $I(A;(\mathbf{c},\{\ell\})) = H(A)$ is computed. The receptor
    vectors $\{\mathbf{v}_u\}$ are then optimized to maximize it.   }
\end{figure*}

{\color{red}
The biological role of an array of sensory receptors is to sense a wide
variety of ligands over a wide range of concentrations. Chemosensitive
systems such as olfaction achieve this through gene expansion, which
generates a diverse repertoire of receptors, each one sensitive to different
chemicals~\cite{valencia-montoya_evolution_2024}. The combined activation of
the array defines a code between the input, namely the identity $\ell$ and
concentration $c$ of the ligands, and the binary response of the $N$ receptors, which we call the activation pattern $A$. 
To study the coding capacity of an array of
chemotactile receptors, we place ourselves upstream of any processing of the signal. Probing a complete olfactome experimentally would require testing
thousands of additional ligands~\cite{mainland_human_2015, munch_door_2016},
so we instead simulate an array of chemotactile receptors, heteromers
included, and compute its coding capacity.
}

Following the standard approach for evaluating the coding capacity of a
receptor array~\cite{fonollosa_quality_2012, del_castillo_convergent_2026}, we
quantify it through the mutual information (MI) between the activation pattern
and the input,
\begin{equation}
    I\bigl(A;(c,\ell)\bigr) = H(A) - H(A \mid c,\ell),
\end{equation}
where $H(A) = -\sum_{A\in\mathbb{A}} p(A)\log p(A)$ is the entropy of the activation pattern
and the mutual information is expressed in bits. We work in the limit of
noiseless receptors, so that $A$ follows deterministically from $(c,\ell)$.
The conditional entropy $H(A \mid c,\ell)$ then vanishes and the mutual
information reduces to the entropy $H(A)$ of the activation pattern under a
given environment.

{\color{blue}
Maximizing $H(A)$ is an instance of the infomax
principle~\cite{bell_information-maximization_1995}, under which information
transfer through a noiseless nonlinear channel is maximized by maximizing the
output entropy. Intuitively, $H(A)$ quantifies how much an organism can measure
about its environment from the activation pattern alone. MI grows with the
number of distinct activation patterns and decreases with the number of
redundant codes, meaning inputs that map to the same pattern and are therefore
indistinguishable to the organism. In the ideal case, all the possible activation patterns are used for a distinct input , which maximizes the entropy up to $H(A) = R$, the number of
receptors.
}

For such an array to operate optimally, it is well established that its
receptors must be promiscuous, each one active for about half of the
ligands, and mutually uncorrelated~\cite{zwicker_receptor_2016,
fonollosa_quality_2012}. Heteromerization greatly expands the combinatorial
family of receptors, but the subunits shared between receptors also make their
responses chemically coupled, and therefore correlated. Whether this coupling
brings additional information or only redundant information is not obvious a
priori. To address this question, we simulate the response of an array to
random sampled ligands.

We first build a physical model of the receptor by extending the
Monod-Wyman-Changeux model of ligand-gated ion channels to heteromers \cite{einav_monod-wyman-changeux_2017}. This
model relates the opening probability of a heteromeric receptor to that of its
homomeric counterparts through the affinities of its constituent subunits for
the ligand, as detailed in the
supplementary material. Because an allosteric channel behaves as a logarithmic
sensor~\cite{olsman_allosteric_2016}, the activation threshold $c^*$ is set on
a logarithmic concentration scale, and the receptor activates once the ligand
concentration exceeds this threshold.

We next design an environment that reproduces random encounters with ligands.
Previous information theoretic models of olfaction do include correlations in
the environment, but only in the concentration and joint occurrence of
ligands~\cite{zwicker_receptor_2016, tesileanu_adaptation_2019}, leaving the
underlying chemistry unstructured and drawing receptor affinities
independently. This is insufficient for our purpose because of a property
specific to heteromers. Receptors that share a subunit respond in a coupled
way, and this coupling, which is the signature of heteromerization, produces
systematic redundancy only when the chemistry of the environment is itself
correlated. We therefore need a realistic morphochemical environment in which
chemically similar ligands occupy nearby regions of chemical space.

We represent the combined morphological and chemical properties of ligands and
binding pockets as vectors embedded in a Euclidean space, denoted
$\mathbf{v}_\ell$ and $\mathbf{v}_u$ respectively and shown in
Fig.~\ref{fig:env}\textbf{A}. Such a Euclidean morphochemical space was first
introduced in the context of immunology~\cite{perelson_theoretical_1979} and
is supported by empirical embeddings of odor space~\cite{lee_principal_2023,
wu_learning_2018}. The activation threshold $c^*$ rises smoothly with the
latent distance $d = \|\mathbf{v}_u - \mathbf{v}_\ell\|$, as shown in
Fig.~\ref{fig:env}\textbf{B}, and saturates between a lower bound
$c^*_\text{min}$ reached at perfect match and an upper bound $c^*_\text{max}$
reached at complete mismatch. This saturating shape is the natural form for a
distribution of binding constants governed by subsite
complementarity~\cite{lancet_probability_1993}, and we implement it as a
Gaussian radial basis kernel (supplementary material).

For each sensing event we draw a mixture of ligands, using a Bernoulli mask
for their presence and lognormal values for their concentrations $\mathbf{c}$,
and pass this mixture through the activation curve to read out a binary
activation pattern $A$. Iterating this Monte Carlo loop yields the empirical
distribution $p(A)$ shown in Fig.~\ref{fig:env}\textbf{C}, from which we
estimate the mutual information. To model how evolution tuned the chemistry of
the receptors to their environment, we optimize the subunit vectors
$\mathbf{v}_u$ so as to maximize the mutual information of the
array~\cite{tesileanu_adaptation_2019, zwicker_receptor_2016}.

Within this framework we can generate random environments of arbitrary
complexity. In fact, we identify a regime of parameters in which our results
become independent of the precise environmental parameters (for details on
the generation of $\mathbf{v}_\ell$, see supplementary material). To
ensure that any difference in mutual information between homomers and
heteromers reflects a property of the array rather than of the environment, we
work in a regime where homomers behave as a perfect
array~\cite{zwicker_receptor_2016}. In this regime the environment does not
limit MI, so that any residual difference is a genuine
characteristic of the receptor array itself. We then compute the mutual information
averaged over many environments, for a growing number of encoding genes and
increasing degrees of heteromerization, as reported in Fig.~\ref{fig:MI}.

Figure~\ref{fig:MI} shows that the mutual information increases with the
heteromerization index, which establishes heteromerization as a viable
strategy for expanding the coding capacity of the array. Two limitations
nonetheless appear. First, the gain decreases as the heteromerization index
grows. Second, expanding the genetic pool to expand receptors' diversity, rather than heteromerization, the MI increases even more.
%at a fixed number of receptors, heteromerization is less effective than gene expansion, as shown in the supplementary material.

{\color{blue}
The first limitation reflects the growing redundancy between receptors as the
heteromerization index increases. In the language of partial information
decomposition~\cite{williams_nonnegative_2010}, the information contributed by
a heteromer that combines many genes becomes increasingly redundant with the
information already carried by heteromers that combine fewer genes. As a
concrete case, the response of a heteromer built from three genes is almost
entirely reconstructible from the set of all receptors built from two genes,
so its unique and synergistic contribution, which is the part that actually
raises $H(A)$, becomes small. The second limitation is expected, because
heteromers are correlated by construction, and a correlated array cannot match
the performance of a perfect array.
}


\begin{figure}
    \centering
    \includegraphics[width=\linewidth]{impact_of_heteromerization.pdf}
    \caption{\label{fig:MI}Mutual information (MI) between the activation pattern and the input, versus the number of encoding genes
    $n_\text{genes}$ for a growing heteromerization index $R/n_\text{genes}$
    ($R$ the number of receptors). $R/n_\text{genes}\approx1$ is gene
    expansion, one gene per receptor. 
    %Solid lines are simulated arrays (bands: spread across environments), dashed lines the perfect homomer array with the same $R$. 
    MI grows with the heteromerization index, but the gain per
    level decays and stays below the perfect reference.
    Error bars grow with $R$ because the phase space of $\mathbb{A}$ scales as $2^R$, making the exact entropy intractable.
    We bracket the Shannon entropy between two bounds : solid lines correspond to an upper-bounded estimate, whereas the dotted lines correspond to a lower-bound, see supplementary material for details.}
    %, so we bracket it between two bounds from pairwise distances between per-event activation distributions~\cite{kolchinsky_estimating_2017}. See supplementary material.}
\end{figure}

\bibliographystyle{apsrev4-2}
\bibliography{smelling}

\end{document}